\documentclass[11pt,a4paper]{article}
\usepackage[utf8]{inputenc}
\usepackage[T1]{fontenc}
\usepackage[english]{babel}
\usepackage{lmodern}
\usepackage[margin=2.3cm]{geometry}
\usepackage{tabularx,booktabs,array,enumitem}
\usepackage[table]{xcolor}
\usepackage{textcomp}
\usepackage{hyperref}

\definecolor{roamly}{HTML}{0E6E68}
\definecolor{pendente}{HTML}{6B7280}
\hypersetup{colorlinks=true,linkcolor=roamly,urlcolor=roamly}
\newcolumntype{L}{>{\raggedright\arraybackslash}X}
\renewcommand{\arraystretch}{1.3}
\setlist{nosep,leftmargin=*}
\setlength{\parindent}{0pt}
\setlength{\parskip}{4pt}

% reference to an interview record or a breakdown (as in the report)
\newcommand{\reg}[1]{{\footnotesize\textcolor{pendente}{(#1)}}}
% use case
\newenvironment{usecase}[2]{%
  \par\medskip\noindent{\large\bfseries\color{roamly}#1: #2}\par\smallskip\small\noindent\ignorespaces}{\par\medskip}
\newcommand{\ucField}[2]{\textbf{#1} & #2 \\}

\begin{document}

\begin{usecase}{UC-04}{Get the group to decide on the proposal}
\textbf{Brief description.} The group coordinator shares a trip proposal with the travel companions, collects each person's possible dates and opinions within a deadline, and records the group's decision so the trip can be booked and paid. While the group decides, the system tells everyone when a price changes, so no one decides on an outdated total.\par\smallskip
\begin{tabularx}{\textwidth}{@{}p{3.4cm}L@{}}
\toprule
\ucField{Primary actor}{Group coordinator}
\ucField{Secondary actors}{Travel companions (give their availability, opinions and votes)}
\ucField{Level}{User goal. The decision is asynchronous: the coordinator starts it in one session and closes it when the group has answered, within a deadline.}
\ucField{Scope}{Roamly app (base version, traveller)}
\ucField{Focus area}{5. Decide with the group or the client}
\ucField{Trigger}{The coordinator has a proposal with the total per person calculated (UC-02) and wants a decision before the price changes. Today, what pushes the organiser to chase the group again is noticing that the price ``went up'' \reg{UE08}; in the system, the automatic price-rise alert takes over that role (extension *a).}
\ucField{Preconditions}{The trip has a proposal with the total per person calculated. The budget per person is recorded (area 1). The coordinator knows who the companions are; they do not need to have opened Roamly before (extension 2a).}
\ucField{Minimal guarantee}{Opinions and dates already given are not lost. No price rise is hidden from the group.}
\ucField{Success guarantee}{The group's decision is recorded and visible to everyone, with the chosen proposal and the total per person. The trip can move on to booking on the supplier's site (area E) and to the group's payments (area 6).}
\bottomrule
\end{tabularx}

\medskip
\textbf{Stakeholders and interests}
\begin{itemize}
  \item \textbf{Coordinator:} get a decision before the price rises \reg{UE08}; not have to chase answers one by one; not be the only one who remembers what was agreed \reg{UE18}.
  \item \textbf{Travel companions:} give an opinion with little effort, without opening spreadsheets \reg{UE17}; understand why an option was cut \reg{UE12}.
  \item \textbf{Family:} know where their child will sleep \reg{UE13}.
\end{itemize}

\textbf{Main success scenario}
\begin{enumerate}
  \item The coordinator shares the proposal with the group and sets a deadline for the decision.
  \item The system shows each companion the proposal, with the total per person and what was cut and why.
  \item Each companion gives their possible dates and their opinion on each option.
  \item The system shows the coordinator everyone's opinions and dates, and who has not answered yet.
  \item The coordinator records the group's decision.
  \item The system saves the decision, shows it to the whole group and lets the trip move on to booking and payments.
\end{enumerate}

\textbf{Extensions}
\begin{itemize}[leftmargin=2.2em]
  \item[1a.] \textit{The proposal's total exceeds the budget per person} \reg{UE06}: the system warns the coordinator before it is shared. The coordinator adjusts the proposal (UC-02) or shares it anyway, with the warning visible to the group.
  \item[2a.] \textit{A companion does not open the proposal or is not in Roamly} \reg{UE17, FM.9, AM.25}: the coordinator sends the group a summary of the proposal that can be seen without opening anything (the ``print'' the group asks for), with the total per person. The coordinator records that companion's answer on their behalf (step 3).
  \item[3a.] \textit{A companion has not answered and the deadline is approaching} \reg{FM.8}: the system reminds them before the deadline. If the deadline passes, the coordinator extends it or decides with the answers available (step 5).
  \item[3b.] \textit{A companion can only travel on other dates} \reg{UE09}: the coordinator redoes the search with the new dates (area 1) and returns to step 1 with the updated proposal.
  \item[3c.] \textit{A companion rejects an option} (``no dorm'', \reg{UE05}): the coordinator swaps the option (area 2) and returns to step 1.
  \item[5b.] \textit{The group is divided on an item or on the dates}: the coordinator changes the items or dates in dispute (areas 1 or 2) and returns to step 1 with the updated proposal.
  \item[5a.] \textit{The group decides in person, around a single phone} \reg{PM.29, SM.22}: the coordinator records what was agreed on the spot, on that phone; the system saves it as the group's decision and shows it to everyone (step 6).
  \item[6a.] \textit{Someone outside the group needs to know the plan} (the mother asks where they will sleep, \reg{UE13}): the coordinator shares the plan with that person without adding them to the group.
  \item[*a.] \textit{The price of an option rises while the group is deciding} \reg{FM.2, SM.17}: the system alerts the group and the coordinator, with the new total per person. If the total now exceeds the budget, continue with 1a.
\end{itemize}

\textbf{Traceability}
\begin{itemize}
  \item \textbf{Consolidated sequence:} activity 6 ``Take the proposal to whoever decides'' (intent: get a decision from the travellers before the price changes); activity 9 ``Bring the group up to date on the costs'' and ``Choose the rest together and record it''; activity 8 ``Tell whoever needs to know''; activity 2 ``Find out the budget, the dates and what can't be given up''.
  \item \textbf{UED, area 5:} ``Show the group the shared idea and proposal'' (step 2), ``Show the group what was cut and why'' (2), ``Enter each friend's possible dates'' (3), ``Give an opinion on each option'' (3), ``Set a deadline for the group to decide'' (1), ``Record what was agreed'' (5, 5a), ``Alert the group when the price rises'' (*a), ``Share the plan with people who aren't going'' (6a). Links to areas 1 (3b), 6 and E (6).
  \item \textbf{Affinity:} ``I wait for the group and meanwhile the price goes up''; ``The proposal I share doesn't reach my friends''.
  \item \textbf{Breakdowns:} FM.2, SM.17 (*a); FM.8 (3a); FM.9, SM.21, AM.25 (step 2: the proposal is seen without opening a spreadsheet; 2a when it still does not reach someone); SM.22, PM.29 (5a).
  \item \textbf{Persona:} Inês.
\end{itemize}

\textbf{Open issues}
\begin{itemize}
  \item When the proposal changes and the use case returns to step 1 (3b, 3c, 5b), it is not known whether opinions and dates given on items that did not change carry over or must be given again. The data does not answer this.
  \item Extensions 3a, 3c and 1a need parts of the system that the UED does not have yet: an automatic reminder to whoever has not answered in area 5 (today there is only one in area 6, for payments), and links from area 5 to area 2 (swap an option) and to area 3 (adjust the proposal). The UED only links area 5 to areas 1 and 6.
  \item In 2a, it is not known how the coordinator records an answer given outside Roamly, or how the group can tell it apart from one the companion gave.
  \item Area 5 also serves the agent's client (A0 shows the proposal to the couple in 5), where the client takes the proposal away to compare online \reg{AE06, FM.5}. That variant is out of this use case's scope and has no use case yet.
\end{itemize}
\end{usecase}

\end{document}
